% ============================================================================
% LEHRERHINWEISE - Sequenz 1a: Rally + Begrüßung
% ============================================================================

\documentclass[11pt,a4paper]{article}

\usepackage[utf8]{inputenc}
\usepackage[T1]{fontenc}
\usepackage[ngerman]{babel}
\usepackage[left=2cm,right=2cm,top=1.8cm,bottom=1.5cm]{geometry}
\usepackage{helvet}
\usepackage{xcolor}
\usepackage{tabularx}
\usepackage{array}
\usepackage{enumitem}
\usepackage{tcolorbox}
\usepackage{fancyhdr}
\usepackage{lastpage}
\usepackage{colortbl}

\renewcommand{\familydefault}{\sfdefault}

\definecolor{spanisch}{HTML}{c41e3a}
\definecolor{grau}{HTML}{666666}
\definecolor{basis}{HTML}{2a7a4b}
\definecolor{standard}{HTML}{2c5aa0}
\definecolor{erweiterung}{HTML}{7b1fa2}

\pagestyle{fancy}
\fancyhf{}
\fancyhead[L]{\small\textcolor{grau}{Spanisch Klasse 7}}
\fancyhead[C]{\small\textcolor{spanisch}{\textbf{LEHRERHINWEISE}}}
\fancyhead[R]{\small\textcolor{grau}{LH-01a}}
\fancyfoot[C]{\small\thepage/\pageref{LastPage}}
\renewcommand{\headrulewidth}{0.5pt}

\newtcolorbox{infobox}[1][]{
    colback=white,
    colframe=spanisch,
    fonttitle=\bfseries\color{white},
    colbacktitle=spanisch,
    title=#1,
    boxrule=1pt,
    arc=0pt,
    left=5pt, right=5pt, top=5pt, bottom=5pt
}

\newcommand{\niveau}[2]{%
    \ifnum#1=1 \textcolor{basis}{\textbf{[B]}} #2%
    \else\ifnum#1=2 \textcolor{standard}{\textbf{[S]}} #2%
    \else \textcolor{erweiterung}{\textbf{[E]}} #2%
    \fi\fi%
}

\begin{document}

% ============================================================================
% KOPFBEREICH
% ============================================================================
\noindent
\begin{tabularx}{\textwidth}{@{}Xr@{}}
    {\Large\textbf{\textcolor{spanisch}{Lehrerhinweise 1a}}} & Sequenz 1: Empezamos \\[0.3em]
    {\small Rally durchs Buch + Begrüßung} & Std. 1-2 (Doppelstunde) \\
\end{tabularx}

\vspace{0.5em}
\hrule
\vspace{1em}

% ============================================================================
% KURZÜBERSICHT
% ============================================================================
\begin{infobox}[Kurzübersicht]
\begin{tabularx}{\textwidth}{@{}lX@{}}
\textbf{Thema:} & Rally durchs Lehrwerk + Begrüßung/Verabschiedung \\
\textbf{Zeit:} & 90 Minuten (Doppelstunde) \\
\textbf{Lehrwerk:} & Qué pasa 1, Empezamos (S. 10-13) \\
\textbf{Materialien:} & AB-01a, LP-01a (digital), Lehrwerk \\
\textbf{Kompetenzen:} & Begrüßen, Verabschieden (A1) \\
\end{tabularx}
\end{infobox}

\vspace{1em}

% ============================================================================
% STUNDENVERLAUF (Kurzform)
% ============================================================================
\section*{Stundenverlauf}

\renewcommand{\arraystretch}{1.3}
\begin{tabularx}{\textwidth}{|c|l|X|l|}
\hline
\rowcolor{spanisch!20}
\textbf{Min} & \textbf{Phase} & \textbf{Inhalt} & \textbf{Sozialform} \\
\hline
\multicolumn{4}{|c|}{\cellcolor{spanisch!10}\textbf{1. Stunde}} \\
\hline
0-5 & Einstieg & Begrüßung auf Spanisch, Motivationsimpuls & Plenum \\
\hline
5-20 & Rally & SuS erkunden Lehrwerk in PA & Partnerarbeit \\
\hline
20-30 & Sicherung & Besprechung Rally-Ergebnisse & Plenum \\
\hline
30-45 & Input & Begrüßungen einführen (Audio) & Plenum \\
\hline
\multicolumn{4}{|c|}{\cellcolor{spanisch!10}\textbf{2. Stunde}} \\
\hline
45-55 & Übung & Zuordnungsaufgabe (AB Aufg. 1) & Einzelarbeit \\
\hline
55-70 & Vertiefung & Tageszeiten + Begrüßung (Aufg. 2) & EA/PA \\
\hline
70-80 & Produktion & Mini-Dialog schreiben (Aufg. 4) & Partnerarbeit \\
\hline
80-90 & Sicherung & Dialoge vorstellen, HA & Plenum \\
\hline
\end{tabularx}

\vspace{1.5em}

% ============================================================================
% DIFFERENZIERUNG
% ============================================================================
\section*{Differenzierung}

\begin{tabularx}{\textwidth}{|l|X|}
\hline
\rowcolor{basis!20}
\niveau{1}{Basis} & Rally: Nur Fragen a-c bearbeiten. Dialog: Lückentext vorgeben. \\
\hline
\rowcolor{standard!20}
\niveau{2}{Standard} & Alle Aufgaben wie im AB vorgesehen. \\
\hline
\rowcolor{erweiterung!20}
\niveau{3}{Erweiterung} & Zusatz: 3 weitere Internationalismen finden. Dialog erweitern. \\
\hline
\end{tabularx}

\vspace{1.5em}

% ============================================================================
% HÄUFIGE FEHLER
% ============================================================================
\section*{Häufige Fehler und Reaktionen}

\begin{enumerate}[leftmargin=*, nosep]
    \item \textbf{Aussprache ``días'' mit deutschem D:}

    \textit{Reaktion:} Auf weiches spanisches D hinweisen, zwischen Zähnen gebildet.

    \item \textbf{Verwechslung ``tardes'' / ``noches'':}

    \textit{Reaktion:} Eselsbrücke: Sonne = días, Dämmerung = tardes, Mond = noches

    \item \textbf{Vergessen der Ausrufezeichen:}

    \textit{Reaktion:} Besonderheit des Spanischen betonen: Satzzeichen am Anfang UND Ende.
\end{enumerate}

\vspace{1.5em}

% ============================================================================
% LÖSUNGEN
% ============================================================================
\section*{Lösungen AB-01a}

\textbf{Merkkasten:}
\begin{itemize}[nosep, leftmargin=*]
    \item Hallo = ¡Hola!
    \item Guten Morgen = Buenos días
    \item Guten Tag = Buenas tardes
    \item Guten Abend = Buenas noches
    \item Tschüss = ¡Adiós!
    \item Bis später = ¡Hasta luego!
    \item Bis morgen = ¡Hasta mañana!
\end{itemize}

\vspace{0.5em}
\textbf{Aufgabe 1:} 1-d, 2-c, 3-b, 4-a

\vspace{0.5em}
\textbf{Aufgabe 2:}
\begin{enumerate}[label=\alph*), nosep, leftmargin=*]
    \item Buenos días
    \item Buenas tardes
    \item Buenas noches
    \item ¡Hola!
\end{enumerate}

\vspace{0.5em}
\textbf{Aufgabe 3:}
\begin{enumerate}[label=\alph*), nosep, leftmargin=*]
    \item 8 Unidades
    \item Seite 196
    \item Hörübung / Audio
    \item z.B. hotel, restaurante, música, chocolate, familia...
\end{enumerate}

\vspace{0.5em}
\textbf{Aufgabe 4 (Beispiel):}\\
A: ¡Hola! Buenos días.\\
B: ¡Hola! Buenos días.\\
A: ¡Hasta luego!\\
B: ¡Adiós! ¡Hasta mañana!

\vspace{1.5em}

% ============================================================================
% HINWEISE
% ============================================================================
\section*{Didaktische Hinweise}

\begin{itemize}[leftmargin=*, nosep]
    \item \textbf{Rally:} Bücher müssen vorhanden sein! Alternativ: Digitale Version nutzen.
    \item \textbf{Audio:} Qué pasa 1 CD Track 1-3 für authentische Aussprache.
    \item \textbf{Lernpfad:} LP-01a als Hausaufgabe oder zur Vertiefung einsetzbar.
    \item \textbf{Erste Stunde:} Bewusst niedrigschwellig für positiven Einstieg.
\end{itemize}

\vspace{1em}

\begin{infobox}[Hausaufgabe]
\begin{itemize}[nosep, leftmargin=*]
    \item Lernpfad LP-01a bearbeiten (digital)
    \item Begrüßungen auswendig lernen
    \item Optional: 5 weitere Internationalismen sammeln
\end{itemize}
\end{infobox}

\end{document}
